\scanpage{53}

% manuscript line(s) 1–2
L'accouplement $\mathcal A\times E\longrightarrow\underline O$ \quad $(x,f)\longmapsto\langle x,f\rangle=f(x)$ est donné par
% manuscript line(s) 3
\[
\boxed{\langle(b,\unsure{r|\gamma}),(\lambda,\unsure{\zeta|\xi})\rangle=\tfrac12(\lambda b+r\cdot\zeta)}
\]
% manuscript line(s) 4
C'est OK car
% manuscript line(s) 5
\[
\underbrace{\lambda_0\underset{\substack{\shortparallel\\r_0\cdot T_0}}{b_0}}_{\lambda_0r_0\cdot T_0}+\underbrace{r_0\cdot\underset{\substack{\deleted{\unsure{r_0T_0}}\shortparallel\\\lambda_0T_0}}{\zeta_0}}_{r_0\lambda_0T_0}=2r_0\lambda_0T_0=0
\]

% manuscript line(s) 6
On a
% manuscript line(s) 7–8
\[
\begin{aligned}\varphi_N\big((\underset{\substack{\in\\\Gamma O}}b,\underset{\substack{\in\\\Gamma L^\vee}}r),(\underset{\substack{\in\\\Gamma\underline O}}{b'},\underset{\substack{\in\\\Gamma L^\vee}}{r'},)\big)&=\tfrac12(bb'\deleted{\unsure{+}}\unsure{-|+}rr'\cdot\delta)\\ &=\langle(\underset{\substack{\in\\\Gamma O}}b,\underset{\substack{\in\\\Gamma L^\vee}}r),(\underset{\substack{\in\\\Gamma O}}{b'},\underbrace{\unsure{-}r'\delta}_{\in\Gamma L})\rangle\end{aligned}
\]
% manuscript line(s) 8
, donc
% manuscript line(s) 9
$f_N:\mathcal A\unsure{\mathbin{-}|\to}E$ est donné par
% manuscript line(s) 10
\[
\boxed{f_N(\underset{\substack{\in\\\Gamma\underline O}}b,\underset{\substack{\in\\\Gamma L^\vee}}r)=(\underset{\substack{\in\\\Gamma\underline O}}b,\underset{\substack{\in\\\Gamma L}}{\deleted{\unsure{+}}\inserted{-}r\delta})}
\]

% manuscript line(s) 11–12
\deleted{Vérifions} \inserted{Utilisons} la formule $Q(f_N(x))=+N(x)\delta$ par la formule précédente, pour $f_N(x)$, $x=(b,r)$, \deleted{\unsure{être}}.
% manuscript line(s) 13–15
\deleted{\unsure{Q}} donc
% manuscript line(s) 13–15
\[
\begin{aligned}Q(\underset{\substack{\shortparallel\\\lambda}}{b},\underbrace{-r\cdot\delta}_{\shortparallel\,\zeta})&=\deleted{\unsure{-}}(b^2-r^2\cdot\delta)/4\ \delta\\ &=\deleted{\unsure{-}}(\lambda^2+r\zeta)/4\ \delta\\ &\text{où}\qquad r=\unsure{-}\zeta\cdot\delta^{-1}\end{aligned}
\]
% manuscript line(s) 16
d'où
% manuscript line(s) 16–17
\[
\begin{aligned}Q(\lambda,\zeta)&=+(\lambda^2\deleted{\unsure{-}}\unsure{+|-}\zeta^2\delta^{-1})/4\ \delta\\ &=\deleted{\unsure{\tfrac14}}(\zeta^2-\lambda^2\delta)/4\end{aligned}
\]
% manuscript line(s) 18
i.e.
% manuscript line(s) 18
\[
\boxed{Q(\lambda,\zeta)=(\zeta^2\deleted{\unsure{+}}-\lambda^2\delta)/4}
\]
% manuscript line(s) 19
Il faut vérifier que $\zeta^2-\lambda^2\delta\equiv0\ (4)$
% manuscript line(s) 19–21
\margin{Or $\zeta\equiv\lambda T(2)$ d'où \deleted{\unsure{la}} $\zeta^2\equiv\lambda^2T^2\ (4)$ donc $\zeta^2-\lambda^2\delta\underset{\unsure{4|d}}\equiv\lambda^2(\underbrace{T^2-\delta}_{4N})$}
% manuscript line(s) 20–21
$=4\lambda^2N\equiv0\ (4)$ OK \deleted{$\zeta_0^2=\lambda_0^2\delta_0$ Or on a $\zeta_0=\lambda_0T_0$ $\zeta_0^2=\lambda_0^2T_0^2=\lambda_0^2\delta_0$ car $T_0^2=\delta_0$ \unsure{OK!}}
% archive sheet number: 52
